\documentclass[9pt]{beamer}

% {{{ PACKAGES
% {{{ ENCODAGE ET POLICE
\usepackage[utf8]{inputenc}
\usepackage[T1]{fontenc}
\usepackage[french]{babel}
\usepackage{eurosym}
% }}}

% {{{ FIGURES
\usepackage{graphicx}
\usepackage{subfig}
% }}}

% {{{ INSERTION DE CODE
\usepackage{moreverb}
% }}}

% {{{ INSERTION DE PDF
\usepackage{pdfpages}
% }}}

% {{{ RÉDUCTION DES MARGES
%\usepackage{fullpage}
%\usepackage[top=2.5cm, bottom=2.5cm, left=3.5cm, right=3.5cm]{geometry}
% }}}

% {{{ MATHÉMATIQUES
\usepackage{amsmath}
\usepackage{amsthm}
\usepackage{bbm}
\usepackage{amssymb}
\usepackage{stmaryrd}
% }}}

% {{{ GRAPHES
\usepackage{tikz}
\usepackage{pgf}
\usetikzlibrary{arrows,automata,shapes}
% }}}

% {{{ RESEAUX DE pETRI
\usetikzlibrary{petri}
% }}}

% {{{ ARBRES
\usetikzlibrary{trees}
% }}}

% {{{ DIAGRAMME DE GANTT
\usepackage{pgfgantt}
% }}}

% {{{ ALGORITHMES
\usepackage{algorithm}
\usepackage[noend]{algpseudocode}
%\usepackage{algorithmic}
%\usepackage{algorithmicx}
% }}}
% }}}

% {{{ NEWTHEOREM
% {{{ THÉORÈMES, lEMMES -- NUMÉROTATION PERSONNALISÉE
\theoremstyle{plain}
\newtheorem{thrm}{Théorème}[section]
\newtheorem{lemm}{Lemme}[subsection]
\newtheorem{corol}{Corollaire}[section]
% }}}

% {{{ THÉORÈMES, lEMMES -- NUMÉROTATION CLASSIQUE
\newtheorem{nthrm}{Théorème}
\newtheorem{nlemm}{Lemme}
\newtheorem{ncorol}{Corollaire}
\newtheorem{npb}{Problème}
% }}}

% {{{ THÉORÈMES, lEMMES -- NON NUMÉROTÉS
\newtheorem*{thrm*}{Théorème}
\newtheorem*{lemm*}{Lemme}
\newtheorem*{corol*}{Corollaire}
% }}}

% {{{ DÉFINITIONS, pROPRIÉTÉS -- NUMÉROTATION PERSONNALISÉE
\theoremstyle{definition}
\newtheorem{df}{Définition}[subsection]
\newtheorem{prop}{Propriété}[subsection]
\newtheorem{conj}{Conjecture}[subsection]
% }}}

% {{{ DÉFINITIONS, pROPRIÉTÉS -- NUMÉROTATION CLASSIQUE
\newtheorem{ndf}{Définition}
\newtheorem{nprop}{Propriété}
\newtheorem{nconj}{Conjecture}
% }}}

% {{{ DÉFINITIONS, pROPRIÉTÉS -- NON NUMÉROTÉS
\newtheorem*{df*}{Définition}
\newtheorem*{prop*}{Propriété}
\newtheorem*{conj*}{Conjecture}
% }}}

% {{{ AUTRES -- NUMÉROTATION PERSONNALISÉE
\theoremstyle{remark}
\newtheorem{nota}{Note}[subsubsection]
\newtheorem{pr}{Proposition}[subsection]
\newtheorem{rmq}{Remarque}[subsubsection]
% }}}

% {{{ AUTRES -- NUMÉROTATION CLASSIQUE
\newtheorem{nnota}{Note}
\newtheorem{npr}{Proposition}
\newtheorem{nrmq}{Remarque}
% }}}

% {{{ AUTRES -- NON NUMÉROTÉS
\newtheorem*{nota*}{Note}
\newtheorem*{pr*}{Proposition}
\newtheorem*{rmq*}{Remarque}
% }}}

% {{{ EXEMPLES
\newtheorem*{ex}{Exemple}
\newtheorem*{exo}{Exercice}
\newtheorem*{rappel}{Rappel}
% }}}
% }}}

% {{{ OPÉRATEURS MATHÉMATIQUES
\DeclareMathOperator{\poly}{\textbf{poly}}
\DeclareMathOperator{\gap}{\mbox{gap}}
\DeclareMathOperator{\Bg}{\mbox{Big}}
\DeclareMathOperator{\Sml}{\mbox{Sml}}
\DeclareMathOperator{\Diff}{\mbox{Diff}}
\DeclareMathOperator{\card}{\mbox{Card}}
\DeclareMathOperator{\Obj}{\mbox{Obj}}
\DeclareMathOperator{\concat}{\bullet}
\DeclareMathOperator{\sig}{\mbox{Sig}}
\DeclareMathOperator{\ver}{\mbox{Ver}}
\DeclareMathOperator{\dist}{\mbox{dist}}
\DeclareMathOperator{\bw}{\mbox{branchwidth}}
\DeclareMathOperator{\tw}{\mbox{tw}}
\DeclareMathOperator{\lm}{\leq_m}
\DeclareMathOperator{\lc}{\leq_c}
\DeclareMathOperator{\lex}{\mbox{lex}}
\DeclareMathOperator{\ord}{\mbox{ord}}
% }}}

% {{{ MACCROS
% {{{ MACROS DIVERSES
\newcommand{\npc}{$NP-$complet}
\newcommand{\npd}{$NP-$difficile}
\newcommand{\la}{largeur arborescente}
\newcommand{\legendre}[2]{\left ( \frac{#1}{#2} \right )}
\newcommand{\wqo}[0]{\emph{well-quasi-ordered}}
\newcommand{\Tau}[0]{\mathcal{T} = (T, \{X_t : t \in T\})}
\newcommand{\Taug}[0]{\mathcal{T}_G = (T, \{X_t : t \in T\})}
\newcommand{\Taustar}[0]{\mathcal{T}^* = (T, \{X_t : t \in T\})}
\newcommand{\Taugstar}[0]{\mathcal{T}_G^* = (T, \{X_t : t \in T\})}
\newcommand{\pdsol}[1]{(X_{#1}^0, X_{#1}^1, X_{#1}^2, M_{#1})}
% }}}
                
% {{{ NOM DE PROBLEMES
\newcommand{\hcycle}{\textsc{HAMILTONIAN CYCLE} }
\newcommand{\phcycle}[1]{\textsc{HAMILTONIAN CYCLE }paramétré par \textsc{#1} }
\newcommand{\vcover}{\textsc{VERTEX COVER} }
\newcommand{\twidth}{\textsc{TREE WIDTH} }
\newcommand{\wiset}{\textsc{WEIGHTED INDEPENDANT SET} }
\newcommand{\fvset}{\textsc{FEEDBACK VERTEX SET} }
\newcommand{\kvcover}{\textsc{K-Vertex Cover} }
\newcommand{\oct}{\textrm{\textsc{Odd Cycle Transversal}} }
\newcommand{\dlp}{\textrm{\textsc{Discrete Logarithm Problem}} }
\newcommand{\ecdlp}{\textrm{\textsc{Elliptic Curve Discrete Logarithm Problem}} }

% }}}

% {{{ MISE EN PAGE

\newcommand{\dfpbi}[3]{
    \begin{npb}[#1]
        \begin{itemize}
            \itemindent 10mm
            \item[\textbf{Données :}] #2
            \item[\textbf{Problème :}] #3
        \end{itemize}
    \end{npb}
}

\newcommand{\dfpb}[3]{
    \begin{npb}[#1]
        \begin{itemize}
            \itemindent 10mm
            \item[\textbf{Données :}] #2
            \item[\textbf{Question :}] #3
        \end{itemize}
    \end{npb}
}

\newcommand{\dfpbp}[4]{
    \begin{npb}[#1]
        \begin{itemize}
            \itemindent 10mm
            \item[\textbf{Données :}] #2
            \item[\textbf{Paramètre :}] #4
            \item[\textbf{Question :}] #3
        \end{itemize}
    \end{npb}
}
% }}}

% }}}
    
% {{{ TIKZ
%%% {{{ STYLE
\tikzstyle{edge}=[
	-,
	draw=black
]

\tikzstyle{sized}=[
    minimum width=9mm
]

\tikzstyle{tvertex}=[
	node distance=15mm,
    inner sep = 1mm
]

\tikzstyle{tvert}=[
	node distance=5mm,
    inner sep=1mm
]

\tikzstyle{tmvertex}=[
    tvertex,
    sized
]

\tikzstyle{tmvert}=[
    tvert,
    sized
]

\tikzstyle{mvertex}=[
    tmvertex,
	circle,
	draw=black,
	fill=black!50
]

\tikzstyle{mvert}=[
    tmvert,
	circle,
	draw=black,
	fill=black!50
]

\tikzstyle{vertex}=[
    tvertex,
	circle,
	draw=black,
	fill=black!50
]

\tikzstyle{vert}=[
    tvert,
	circle,
	draw=black,
	fill=black!50
]

\tikzstyle{treedec}=[
	vertex,
	node distance=15mm,
	draw=red,
	fill=red!50
]

\tikzstyle{treeedge}=[
	edge,
	draw=red
]

\tikzstyle{giedge}=[
	edge,
	draw=blue
]

\tikzstyle{bounding}=[
	thick, %ultra
	%dash pattern=on 1mm off 1mm,
	densely dotted,
	fill=blue!10,
	draw=blue!70
]

\tikzstyle{patate}=[
    rounded corners=5pt,
    draw=black,
    fill=black!50,
    text width=12mm,
    text centered
]

\tikzstyle{patatoid}=[
    ellipse,
    node distance=10mm,
    draw=black
]
%%% }}}

% {{{ DESSINS
\newcommand{\base}[4]{
    \draw[->] (0,0) to (0, #1);
    \draw[->] (0,0) to (#2, 0);
    
    \node at (0, #1+2) {#3};
    \node at (#2 + 2, -2) {#4};
}
% }}}
%}}}



\usetheme[compress]{Dresden2}
\usecolortheme{beaver}
\usecolortheme{seahorse}
%%\usecolortheme{spruce}
%
\setbeamertemplate{blocks}[rounded][shadow=true]
%\setbeamertemplate{navigation symbols}{} 
%%\setbeamercolor{block body alerted}{bg=red!30}
%%\setbeamercolor{block title alerted}{fg=white,bg=red}

\title[Compression par matroides]{Compression via Matroids : A Randomized Polynomial Kernel for Odd
Cycle Transversal}
\author[G. Duvillié]{Guillerme Duvillié}
\institute{Université Montpellier 2}
%\logo{\includegraphics[scale=0.03]{um2.png}}

\AtBeginSection[]{
    \begin{frame}{Plan}
        \tableofcontents[currentsection]
    \end{frame}
}

\floatname{algorithm}{Procedure}
\renewcommand{\algorithmicrequire}{\textbf{Input:}}
\renewcommand{\algorithmicensure}{\textbf{Output:}}

\begin{document}

\begin{frame}
    \titlepage
\end{frame}

\addtocounter{framenumber}{-1}
\addtobeamertemplate{footline}{\insertframenumber/\inserttotalframenumber}

\section{Introduction}
\subsection*{a}

% {{{ First Frame
\begin{frame}
    \frametitle{\rmfamily{\scshape{Odd Cycle Transversal}} }
    \dfpbp{\rmfamily{\scshape{Odd Cycle Transversal}} }{Un graphe $G=(V,E)$, un entier $k$}{Est-il possible
    de transformer $G$ en un graphe biparti en supprimant au plus $k$ sommets de $G$?}{$k$}

    \vfill

    \pause

    \begin{columns}
        \begin{column}[c]{0.35\linewidth}
            \begin{figure}
                \begin{center}
                    \begin{tikzpicture}
                        \node[vertex] (1) at (0,2) {$1$};
                        \node[vertex] (2) at (2,2) {$2$};
                        \node[vertex] (3) at (3,1) {$3$};
                        \node[vertex] (4) at (2,0) {$4$};
                        \node[vertex] (5) at (1,1) {$5$};
                        \node[vertex] (6) at (0,0) {$6$};

                        \draw[-] (1) to (2);
                        \draw[-] (1) to (5);
                        \draw[-] (1) to (6);

                        \draw[-] (2) to (3);
                        \draw[-] (2) to (4);

                        \draw[-] (3) to (4);

                        \draw[-] (4) to (5);
                        \draw[-] (4) to (6);

                        \draw[-] (5) to (6);
                    \end{tikzpicture}
                \end{center}
                \caption{Exemple}
            \end{figure}
        \end{column}
        \hfill
        \begin{column}[c]{0.45\linewidth}
        \end{column}
    \end{columns}
\end{frame}

\addtocounter{framenumber}{-1}

\begin{frame}
    \frametitle{\rmfamily{\scshape{Odd Cycle Transversal}} }
    \dfpbp{\rmfamily{\scshape{Odd Cycle Transversal}} }{Un graphe $G=(V,E)$, un entier $k$}{Est il possible
    de transformer $G$ en un graphe biparti en supprimant au plus $k$ sommets de $G$?}{$k$}

    \vfill

    \begin{columns}
        \begin{column}[c]{0.35\linewidth}
            \begin{figure}
                \begin{center}
                    \begin{tikzpicture}
                        \node[vertex, fill=red!50] (1) at (0,2) {$1$};
                        \node[vertex, fill=green!50] (2) at (2,2) {$2$};
                        \node[vertex, fill=red!50] (3) at (3,1) {$3$};
                        \node[vertex, fill=blue!50] (4) at (2,0) {$4$};
                        \node[vertex, fill=blue!50] (5) at (1,1) {$5$};
                        \node[vertex, fill=green!50] (6) at (0,0) {$6$};

                        \draw[-] (1) to (2);
                        \draw[-, dotted, blue] (1) to (5);
                        \draw[-] (1) to (6);

                        \draw[-] (2) to (3);
                        \draw[-, dotted, blue] (2) to (4);

                        \draw[-, dotted, blue] (3) to (4);

                        \draw[-, dotted, blue] (4) to (5);
                        \draw[-, dotted, blue] (4) to (6);

                        \draw[-, dotted, blue] (5) to (6);
                    \end{tikzpicture}
                \end{center}
                \caption{Exemple}
            \end{figure}
        \end{column}
        \hfill

        \pause

        \begin{column}[l]{0.45\linewidth}
            \begin{thrm}
                \oct est \npc
            \end{thrm}
        \end{column}
    \end{columns}
\end{frame}
%}}}
% {{{ Second Frame
\begin{frame}
    \frametitle{Fixed Parameter Tractable ?}

    \begin{thrm}[Reed, Smith, Vetta (2003)]
        \oct est FPT: \\
        $\exists$ un algorithme résolvant le problème en $O(4^k mn)$
    \end{thrm}

    \pause
    \vfill

    \begin{exampleblock}{Question}
        \oct possède-t-il un noyau?
    \end{exampleblock}

    \pause
    \vfill

    \begin{alertblock}{Réponse}
        Oui, un noyau randomisé polynomial
    \end{alertblock}
\end{frame}
% }}}

\section[Algorithme]{Principe de l'algorithme FPT}
\subsection*{a}

% {{{ Intro Algo
\begin{frame}
    \frametitle{Un algorithme FPT pour \oct}

    \oct$(k)$
    \begin{itemize} 
        \item[] Un graphe $G=(V,E)$, un entier $k$
        \item[] Un ensemble de sommets $Y$, tel que $|Y| \leq k$ et $Y$ intersecte tous les
            cycles impairs de $G$, si un tel ensemble existe. Faux sinon.
        \item[]
        \item Choisir un sommet $v \in V$
        \item Appeler \oct$(k)$ sur $G - v$
        \item Si $G-v$ n'a pas de solution $G$ non plus
        \item Si $Y \leq k$ retourner $Y$
        \item Si $Y \geq k$ déterminer (s'il existe) un OCT $Y'$ tq $|Y'| \leq k$.
    \end{itemize}

    \pause
    \vfill

    \begin{alertblock}{Question}
        Comment déterminer $Y'$ à partir de $Y$ ?
    \end{alertblock}
\end{frame}
% }}}
% {{{ Transformation Algo
\begin{frame}
    \frametitle{Algorithme : Transformation de $G$ en $G'$}
    Soit un graphe $G$, un \oct $X$ et une $2$-partition $(S_1, S_2)$ de $G - X$
    \vfill

    \begin{figure}
        \begin{center}
            \begin{tikzpicture}
                \node[vertex, fill=red!50] (1) at (0,2) {$1$};
                \node[vertex, fill=green!50] (2) at (2,2) {$2$};
                \node[vertex, fill=red!50] (3) at (3,1) {$3$};
                \node[vertex, fill=blue!50] (4) at (2,0) {$4$};
                \node[vertex, fill=blue!50] (5) at (1,1) {$5$};
                \node[vertex, fill=green!50] (6) at (0,0) {$6$};

                \draw[-] (1) to (2);
                \draw[-, dotted, blue] (1) to (5);
                \draw[-] (1) to (6);

                \draw[-] (2) to (3);
                \draw[-, dotted, blue] (2) to (4);

                \draw[-, dotted, blue] (3) to (4);

                \draw[-, dotted, blue] (4) to (5);
                \draw[-, dotted, blue] (4) to (6);

                \draw[-, dotted, blue] (5) to (6);
            \end{tikzpicture}
        \end{center}
        \caption{Le graphe $G$}
    \end{figure}
\end{frame}

\addtocounter{framenumber}{-1}

\begin{frame}
    \frametitle{Algorithme : Transformation de $G$ en $G'$}
    Les sommets et arêtes de $G - X$ appartiennent à $G'$
    \vfill

    \begin{columns}
        \begin{column}{0.40\linewidth}
            \begin{figure}
                \begin{center}
                    \begin{tikzpicture}
                        \node[vertex, fill=red!50] (1) at (0,2) {$1$};
                        \node[vertex, fill=green!50] (2) at (2,2) {$2$};
                        \node[vertex, fill=red!50] (3) at (3,1) {$3$};
                        \node[vertex, fill=blue!50] (4) at (2,0) {$4$};
                        \node[vertex, fill=blue!50] (5) at (1,1) {$5$};
                        \node[vertex, fill=green!50] (6) at (0,0) {$6$};

                        \draw[-] (1) to (2);
                        \draw[-, dotted, blue] (1) to (5);
                        \draw[-] (1) to (6);

                        \draw[-] (2) to (3);
                        \draw[-, dotted, blue] (2) to (4);

                        \draw[-, dotted, blue] (3) to (4);

                        \draw[-, dotted, blue] (4) to (5);
                        \draw[-, dotted, blue] (4) to (6);

                        \draw[-, dotted, blue] (5) to (6);
                    \end{tikzpicture}
                \end{center}
                \caption{Le graphe $G$}
            \end{figure}
        \end{column}
        \begin{column}{0.05\linewidth}
            $\Longrightarrow$
        \end{column}
        \begin{column}{0.40\linewidth}
            \begin{figure}
                \begin{center}
                    \begin{tikzpicture}
                        \node[vertex, fill=red!50] (1) at (0,3) {$1$};
                        \node[vertex, fill=red!50] (3) at (0,1) {$3$};
                        \node[vertex, fill=green!50] (2) at (3,1) {$2$};
                        \node[vertex, fill=green!50] (6) at (3,3) {$6$};

                        \draw[-] (1) to (2);
                        \draw[-] (1) to (6);

                        \draw[-] (2) to (3);
                    \end{tikzpicture}
                \end{center}
                \caption{Le graphe $G'$}
            \end{figure}
        \end{column}
    \end{columns}
\end{frame}

\addtocounter{framenumber}{-1}

\begin{frame}
    \frametitle{Algorithme : Transformation de $G$ en $G'$}
    Si $v \in X$ on crée $v_1$ et $v_2$ dans $G'$.
    \vfill

    \begin{columns}
        \begin{column}{0.40\linewidth}
            \begin{figure}
                \begin{center}
                    \begin{tikzpicture}
                        \node[vertex, fill=red!50] (1) at (0,2) {$1$};
                        \node[vertex, fill=green!50] (2) at (2,2) {$2$};
                        \node[vertex, fill=red!50] (3) at (3,1) {$3$};
                        \node[vertex, fill=blue!50] (4) at (2,0) {$4$};
                        \node[vertex, fill=blue!50] (5) at (1,1) {$5$};
                        \node[vertex, fill=green!50] (6) at (0,0) {$6$};

                        \draw[-] (1) to (2);
                        \draw[-, dotted, blue] (1) to (5);
                        \draw[-] (1) to (6);

                        \draw[-] (2) to (3);
                        \draw[-, dotted, blue] (2) to (4);

                        \draw[-, dotted, blue] (3) to (4);

                        \draw[-, dotted, blue] (4) to (5);
                        \draw[-, dotted, blue] (4) to (6);

                        \draw[-, dotted, blue] (5) to (6);
                    \end{tikzpicture}
                \end{center}
                \caption{Le graphe $G$}
            \end{figure}
        \end{column}
        \begin{column}{0.05\linewidth}
            $\Longrightarrow$
        \end{column}
        \begin{column}{0.40\linewidth}
            \begin{figure}
                \begin{center}
                    \begin{tikzpicture}
                        \node[vertex, fill=red!50] (1) at (0,3) {$1$};
                        \node[vertex, fill=red!50] (3) at (0,1) {$3$};
                        \node[vertex, fill=green!50] (2) at (3,1) {$2$};
                        \node[vertex, fill=green!50] (6) at (3,3) {$6$};

                        \node[vertex, fill=blue!50] (51) at (1,4) {$5_1$};
                        \node[vertex, fill=blue!50] (52) at (2,4) {$5_2$};
                        \node[vertex, fill=blue!50] (41) at (1,0) {$4_1$};
                        \node[vertex, fill=blue!50] (42) at (2,0) {$4_2$};

                        \draw[-] (1) to (2);
                        \draw[-] (1) to (6);

                        \draw[-] (2) to (3);
                    \end{tikzpicture}
                \end{center}
                \caption{Le graphe $G'$}
            \end{figure}
        \end{column}
    \end{columns}
\end{frame}

\addtocounter{framenumber}{-1}

\begin{frame}
    \frametitle{Algorithme : Transformation de $G$ en $G'$}
    \vfill

    \begin{columns}
        \begin{column}{0.40\linewidth}
            \begin{figure}
                \begin{center}
                    \begin{tikzpicture}
                        \node[vertex, fill=red!50] (1) at (0,2) {$1$};
                        \node[vertex, fill=green!50] (2) at (2,2) {$2$};
                        \node[vertex, fill=red!50] (3) at (3,1) {$3$};
                        \node[vertex, fill=blue!50] (4) at (2,0) {$4$};
                        \node[vertex, fill=blue!50] (5) at (1,1) {$5$};
                        \node[vertex, fill=green!50] (6) at (0,0) {$6$};

                        \draw[-] (1) to (2);
                        \draw[-, dotted, blue] (1) to (5);
                        \draw[-] (1) to (6);

                        \draw[-] (2) to (3);
                        \draw[-, dotted, blue] (2) to (4);

                        \draw[-, dotted, blue] (3) to (4);

                        \draw[-, dotted, blue] (4) to (5);
                        \draw[-, dotted, blue] (4) to (6);

                        \draw[-, dotted, blue] (5) to (6);
                    \end{tikzpicture}
                \end{center}
                \caption{Le graphe $G$}
            \end{figure}
        \end{column}
        \begin{column}{0.05\linewidth}
            $\Longrightarrow$
        \end{column}
        \begin{column}{0.40\linewidth}
            \begin{figure}
                \begin{center}
                    \begin{tikzpicture}
                        \node[vertex, fill=red!50] (1) at (0,3) {$1$};
                        \node[vertex, fill=red!50] (3) at (0,1) {$3$};
                        \node[vertex, fill=green!50] (2) at (3,1) {$2$};
                        \node[vertex, fill=green!50] (6) at (3,3) {$6$};

                        \node[vertex, fill=blue!50] (51) at (1,4) {$5_1$};
                        \node[vertex, fill=blue!50] (52) at (2,4) {$5_2$};
                        \node[vertex, fill=blue!50] (41) at (1,0) {$4_1$};
                        \node[vertex, fill=blue!50] (42) at (2,0) {$4_2$};

                        \draw[-] (1) to (2);
                        \draw[-] (1) to (6);
                        \draw[-] (2) to (3);

                        \draw[-, blue] (1) to (52);
                        \draw[-, blue] (6) to (51);
                        \draw[-, blue] (6) to (41);
                        \draw[-, blue] (2) to (41);
                        \draw[-, blue] (3) to (42);
                        \draw[-, blue] (41) to (52);
                        \draw[-, blue] (42) to (51);
                    \end{tikzpicture}
                \end{center}
                \caption{Le graphe $G'$}
            \end{figure}
        \end{column}
    \end{columns}

    \pause
    \vfill

    \begin{df}
        Étant donné un ensemble $Y \subseteq X$, une partition $(Y_A, Y_B)$ de $Y' = \{y_1, y2 : y
        \in Y \}$ est valide si $\forall y \in Y, |\{y_1, y_2\}| = 1$.
    \end{df}
\end{frame}
% }}}
% {{{ Propriétés Algo
\begin{frame}
    \frametitle{Algorithme : Quelques propriétés}

    \begin{rappel}
        Soient deux ensembles $S, T \subseteq V$, $T$ est lié à $S$ s'il existe $|T|$ chemins
        disjoints de $S$ vers $T$
    \end{rappel}

    \pause
    \vfill

    \begin{rmq}
        $S$ et $T$ ne sont pas nécessairement disjoints et un sommet est lié à lui même par un
        chemin de longueur $0$.
    \end{rmq}

    \pause
    \vfill

    \begin{thrm}
        Un \oct $X$ est de taille minimum ssi pour n'importe quelle partition valide $(Y_A, Y_B)$ de
        n'importe quel sous ensemble $Y$ de $X$, $Y_B$ est lié à $Y_A$ dans $G' - \{x_1, x_2 : x \in
        X - Y \}$
    \end{thrm}
\end{frame}
%}}}
% {{{ Exemple Algo
\begin{frame}
    \frametitle{Algorithme : Illustration du théorème}
    \begin{columns}
        \begin{column}{0.55\linewidth}
            \begin{figure}
                \begin{center}
                    \begin{tikzpicture}[scale=1.2]
                        \tikzset{vertex/.style={circle, draw=black, minimum width=8mm}};
                        \node[vertex, fill=red!50] (1) at (0,3) {$1$};
                        \node[vertex, fill=red!50] (3) at (0,1) {$3$};
                        \node[vertex, fill=green!50] (2) at (3,1) {$2$};
                        \node[vertex, fill=green!50] (6) at (3,3) {$6$};

                        \node[vertex, fill=blue!50] (51) at (1,4) {$5_1$};
                        \node[vertex, fill=blue!50] (52) at (2,4) {$5_2$};
                        \node[vertex, fill=blue!50] (41) at (1,0) {$4_1$};
                        \node[vertex, fill=blue!50] (42) at (2,0) {$4_2$};

                        \draw[-] (1) to (2);
                        \draw[-] (1) to (6);
                        \draw[-] (2) to (3);

                        \draw[-] (1) to (52);
                        \draw[-] (6) to (51);
                        \draw[-] (6) to (41);
                        \draw[-] (2) to (41);
                        \draw[-] (3) to (42);
                        \draw[-] (41) to (52);
                        \draw[-] (42) to (51);
                    \end{tikzpicture}
                \end{center}
                \caption{Le graphe $G'$}
            \end{figure}
        \end{column}
        %\vfill
        \begin{column}{0.35\linewidth}
        \end{column}
    \end{columns}
\end{frame}

\addtocounter{framenumber}{-1}

\begin{frame}
    \frametitle{Algorithme : Illustration du théorème}
    \begin{columns}
        \begin{column}{0.5\linewidth}
            \begin{figure}
                \begin{center}
                    \begin{tikzpicture}[scale=1.2]
                        \tikzset{vertex/.style={circle, draw=black, minimum width=8mm}};
                        \node[vertex, fill=red!50] (1) at (0,3) {$1$};
                        \node[vertex, fill=red!50] (3) at (0,1) {$3$};
                        \node[vertex, fill=green!50] (2) at (3,1) {$2$};
                        \node[vertex, fill=green!50] (6) at (3,3) {$6$};

                        \node at (1,4) {};
                        \node at (2,4) {};

                        %\node[vertex, fill=blue!50] (51) at (1,4) {$5_1$};
                        %\node[vertex, fill=blue!50] (52) at (2,4) {$5_2$};
                        \node[vertex, fill=blue!50] (41) at (1,0) {$4_1$};
                        \node[vertex, fill=blue!50] (42) at (2,0) {$4_2$};

                        \draw[-] (1) to (2);
                        \draw[-] (1) to (6);
                        \draw[-] (6) to (41);
                        \draw[-, red] (2) to (3);

                        %\draw[-, blue] (1) to (52);
                        %\draw[-, blue] (6) to (51);
                        \draw[-, red] (2) to (41);
                        \draw[-, red] (3) to (42);
                        %\draw[-, blue] (41) to (52);
                        %\draw[-, blue] (42) to (51);
                    \end{tikzpicture}
                \end{center}
                \caption{Le graphe $G'$}
            \end{figure}
        \end{column}
        \begin{column}{0.4\linewidth}
            \begin{itemize}
                \item $Y_A = \{4_1\},\ Y_B = \{4_2\}$
            \end{itemize}
        \end{column}
    \end{columns}
\end{frame}

\addtocounter{framenumber}{-1}

\begin{frame}
    \frametitle{Algorithme : Illustration du théorème}
    \begin{columns}
        \begin{column}{0.5\linewidth}
            \begin{figure}
                \begin{center}
                    \begin{tikzpicture}[scale=1.2]
                        \tikzset{vertex/.style={circle, draw=black, minimum width=8mm}};
                        \node[vertex, fill=red!50] (1) at (0,3) {$1$};
                        \node[vertex, fill=red!50] (3) at (0,1) {$3$};
                        \node[vertex, fill=green!50] (2) at (3,1) {$2$};
                        \node[vertex, fill=green!50] (6) at (3,3) {$6$};

                        \node at (1,4) {};
                        \node at (2,4) {};
                        \node at (1,0) {};
                        \node at (2,0) {};

                        \node[vertex, fill=blue!50] (51) at (1,4) {$5_1$};
                        \node[vertex, fill=blue!50] (52) at (2,4) {$5_2$};
                        %\node[vertex, fill=blue!50] (41) at (1,0) {$4_1$};
                        %\node[vertex, fill=blue!50] (42) at (2,0) {$4_2$};

                        \draw[-] (1) to (2);
                        \draw[-, red] (1) to (6);
                        \draw[-] (2) to (3);

                        \draw[-, red] (1) to (52);
                        \draw[-, red] (6) to (51);
                        %\draw[-, red] (2) to (41);
                        %\draw[-, red] (3) to (42);
                        %\draw[-] (6) to (41);
                        %\draw[-, blue] (41) to (52);
                        %\draw[-, blue] (42) to (51);
                    \end{tikzpicture}
                \end{center}
                \caption{Le graphe $G'$}
            \end{figure}
        \end{column}
        \begin{column}{0.4\linewidth}
            \begin{itemize}
                \item $Y_A = \{4_1\},\ Y_B = \{4_2\}$
                \item $Y_A = \{5_1\},\ Y_B = \{5_2\}$
            \end{itemize}
        \end{column}
    \end{columns}
\end{frame}

\addtocounter{framenumber}{-1}

\begin{frame}
    \frametitle{Algorithme : Illustration du théorème}
    \begin{columns}
        \begin{column}{0.5\linewidth}
            \begin{figure}
                \begin{center}
                    \begin{tikzpicture}[scale=1.2]
                        \tikzset{vertex/.style={circle, draw=black, minimum width=8mm}};
                        \node[vertex, fill=red!50] (1) at (0,3) {$1$};
                        \node[vertex, fill=red!50] (3) at (0,1) {$3$};
                        \node[vertex, fill=green!50] (2) at (3,1) {$2$};
                        \node[vertex, fill=green!50] (6) at (3,3) {$6$};

                        \node at (1,4) {};
                        \node at (2,4) {};
                        \node at (1,0) {};
                        \node at (2,0) {};

                        \node[vertex, fill=blue!50] (51) at (1,4) {$5_1$};
                        \node[vertex, fill=blue!50] (52) at (2,4) {$5_2$};
                        \node[vertex, fill=blue!50] (41) at (1,0) {$4_1$};
                        \node[vertex, fill=blue!50] (42) at (2,0) {$4_2$};

                        \draw[-] (1) to (2);
                        \draw[-] (1) to (6);
                        \draw[-] (2) to (3);

                        \draw[-] (1) to (52);
                        \draw[-] (6) to (51);
                        \draw[-] (2) to (41);
                        \draw[-] (3) to (42);
                        \draw[-] (6) to (41);
                        \draw[-, red] (41) to (52);
                        \draw[-, purple] (42) to (51);
                    \end{tikzpicture}
                \end{center}
                \caption{Le graphe $G'$}
            \end{figure}
        \end{column}
        \begin{column}{0.4\linewidth}
            \begin{itemize}
                \item $Y_A = \{4_1\},\ Y_B = \{4_2\}$
                \item $Y_A = \{5_1\},\ Y_B = \{5_2\}$
                \item $Y_A = \{4_1, 5_1\},\ Y_B = \{4_2, 5_2\}$
            \end{itemize}
        \end{column}
    \end{columns}
\end{frame}

\addtocounter{framenumber}{-1}

\begin{frame}
    \frametitle{Algorithme : Illustration du théorème}
    \begin{columns}
        \begin{column}{0.5\linewidth}
            \begin{figure}
                \begin{center}
                    \begin{tikzpicture}[scale=1.2]
                        \tikzset{vertex/.style={circle, draw=black, minimum width=8mm}};
                        \node[vertex, fill=red!50] (1) at (0,3) {$1$};
                        \node[vertex, fill=red!50] (3) at (0,1) {$3$};
                        \node[vertex, fill=green!50] (2) at (3,1) {$2$};
                        \node[vertex, fill=green!50] (6) at (3,3) {$6$};

                        \node at (1,4) {};
                        \node at (2,4) {};
                        \node at (1,0) {};
                        \node at (2,0) {};

                        \node[vertex, fill=blue!50] (51) at (1,4) {$5_1$};
                        \node[vertex, fill=blue!50] (52) at (2,4) {$5_2$};
                        \node[vertex, fill=blue!50] (41) at (1,0) {$4_1$};
                        \node[vertex, fill=blue!50] (42) at (2,0) {$4_2$};

                        \draw[-, purple] (1) to (2);
                        \draw[-] (1) to (6);
                        \draw[-, purple] (2) to (3);

                        \draw[-, purple] (1) to (52);
                        \draw[-, red] (6) to (51);
                        \draw[-] (2) to (41);
                        \draw[-, purple] (3) to (42);
                        \draw[-, red] (6) to (41);
                        \draw[-] (41) to (52);
                        \draw[-] (42) to (51);
                    \end{tikzpicture}
                \end{center}
                \caption{Le graphe $G'$}
            \end{figure}
        \end{column}
        \begin{column}{0.4\linewidth}
            \begin{itemize}
                \item $Y_A = \{4_1\},\ Y_B = \{4_2\}$
                \item $Y_A = \{5_1\},\ Y_B = \{5_2\}$
                \item $Y_A = \{4_1, 5_1\},\ Y_B = \{4_2, 5_2\}$
                \item $Y_A = \{4_1, 5_2\},\ Y_B = \{4_2, 5_1\}$
            \end{itemize}
        \end{column}
    \end{columns}
\end{frame}
%}}}
%{{{ Conclu Algo
\begin{frame}
    \frametitle{Algorithme : Conclusion}
    \begin{rmq}
        La recherche de $k$ chemins disjoints se ramène à un calcul de flot maximum (et donc de
        coupe minimum).
    \end{rmq}

    \vfill
    \pause

    \begin{alertblock}{Question}
        Pourquoi vous avoir présenté tout ça?
    \end{alertblock}

    \vfill
    \pause

    \begin{exampleblock}{Réponse}
        La représentation de l'ensemble des problèmes de coupe résolus par l'algorithme peut être
        polynomiale.
    \end{exampleblock}
\end{frame}
%}}}

\section{Matroides}
\subsection*{a}
% {{{ Matroide : Intro
\begin{frame}
    \frametitle{Un Matroide? Kézako?}
    \begin{columns}
        \begin{column}{0.40\linewidth}
            $G = (V, E)$ :
            \begin{tikzpicture}
                \node[vertex] (a) at (1, 0) {};
                \node[vertex] (b) at (0, 1) {};
                \node[vertex] (c) at (2, 1) {};
                \node[vertex] (d) at (1, 2) {};

                \draw (a) --node[below left] {$1$} (b) --node[above left] {$2$} (d) 
                --node[above right] {$3$} (c) -- node[below right] {$4$} (a) -- (b) --node[above]
                {$5$} (c);
            \end{tikzpicture}
        \end{column}
        \pause
        \begin{column}{0.40\linewidth}
            \[
                A_{\mathbb{F}} = 
                \left [
                \begin{array}{ccccc}
                    1 & 2 & 3 & 4 & 5 \\ \hline
                    1 & 0 & 1 & 0 & 1 \\
                    0 & 1 & 0 & 1 & 1 \\
                    0 & 1 & 1 & 0 & 0 \\
                \end{array}
                \right ]
            \]
        \end{column}
    \end{columns}
    \pause
    \begin{center}
        \begin{tikzpicture}
            \node[inner sep=0mm] (a) at (2, 0.0) {$1$};
            \node[inner sep=0mm] (b) at (1, 0.4) {$4$};
            \node[inner sep=0mm] (c) at (0, 0.8) {$5$};
            \node[inner sep=0mm] (d) at (1, 1.2) {$2$};
            \node[inner sep=0mm] (e) at (2, 1.6) {$3$};

            \draw (a) -- (b) -- (c) -- (d) -- (e);
        \end{tikzpicture}
    \end{center}

    \vfill
    \pause

    \begin{alertblock}{Intuition}
        Un matroide permet de représenter le concept d'indépendance linéaire.
    \end{alertblock}
\end{frame}
% }}}
% {{{ Matroide : Définition formelle
\begin{frame}
    \frametitle{De manière plus formelle}
    \begin{df}[Matroide]
        Un matroide $M$ est un couple $(E, \mathcal{I})$, où $E$ est un ensemble fini et
        $\mathcal{I}$ est une famille de sous-ensembles de $E$ vérifiant les conditions suivantes :
        \begin{enumerate}
            \item $\emptyset \in \mathcal{I}$
            \item Si $I_1 \in \mathcal{I}$ et $I_2 \subseteq I_1$, alors $I_2 \in \mathcal{I}$
            \item Si $I_1, I_2 \in \mathcal{I}$ et $|I_1| < |I_2|$ alors $\exists e \in I_2
                \backslash I_1$ tel que $I_1 \cup \{e\} \in \mathcal{I}$
        \end{enumerate}
    \end{df}

    \vfill
    \pause

    \begin{prop}
        Soit $A$ une matrice dont les coefficients sont à valeur dans un corps fini $\mathbb{F}$ et
        $E$ les ensembles de colonnes de $A$. Soit $\mathcal{I}$ l'ensemble des ensembles $X
        \subseteq E$ des colonnes linéairement indépendantes est un matroide.
    \end{prop}

    \vfill
    \pause

    \begin{df}
        Un matroide est dit représentable s'il existe une matrice dont les coefficients sont à
        valeur dans $\mathbb{F}$ qui le représente.
    \end{df}

\end{frame}
%}}}
%{{{ Résultats sur les matroides
\begin{frame}
    \begin{thrm}[Marx, 2007]
        Soit $D = (V, A)$ un graphe orienté et soit $S \subseteq V$, les sous-ensembles $T \subseteq
        V$ liés à $S$ forment les indépendants d'un matroide. Une représentation de ce matroide peut
        être obtenu à l'aide d'un algorithme randomisé s'exécutant en temps polynomial avec une
        erreur unilatérale.
    \end{thrm}

    \pause
    \vfill

    \begin{corol}
        Étant donnés un graphe orienté $D = (V, A)$, un réel $\epsilon$ et $S$ et $T$ deux
        sous-ensembles de $V$, soit $M$ le matroide formé par les sous-ensembles $T$ liés à $S$. Une
        représentation de $M$ sous forme d'une matrice $|S| \times |T|$ à valeur dans $\mathbb{Q}$
        peut être calculée à l'aide d'un algorithme randomisé s'exécutant en temps polynomial avec
        une erreur unilatérale au plus $\epsilon$.
    \end{corol}

    \pause
    \vfill

    \begin{rmq}
        Les coefficients de la matrice ont une longueur : \[
            O(\min(|T|, |S| \log |T|) + \log  \left ( \frac{1}{\epsilon} \right ) + \log |V| )
        \]
    \end{rmq}
\end{frame}
%}}}
%{{{ Matroide : Intuition
\begin{frame}
    \begin{alertblock}{Problème}
        Ce matroide ne permet de représenter que les sous-ensembles liés à un seul ensemble $S$
    \end{alertblock}

    \pause
    \vfill

    \begin{exampleblock}{Généralisation}
        Soit $G = (V, E)$ un graphe non orienté, soit $X \subseteq V$ un ensemble de sommets, et
        soit $X' = \{v' : v \in X\}$ un ensemble de nouveaux sommets. Il existe une construction
        en temps polynomial d'un graphe orienté $D = (V \cup X', A)$ tel que $I \subseteq X \cup X'$
        est un indépendant du matroide $(D, X')$ ssi $T$ est lié à $S$ dans $G - R$ avec :
        \begin{itemize}
            \item $S = \{v \in X : v, v' \not \in I\}$
            \item $T = \{v \in X : v, v' \in I\}$
            \item $R = \{v \in X : v \in I \mbox{ et } v' \not \in I\}$
        \end{itemize}
    \end{exampleblock}

    \pause
    \vfill

    \begin{rmq}
        Il existe un algorithme randomisé s'exécutant en temps polynomial calculant une matrice de
        taille $|X| \times 2|X|$ dont les coefficients sont des entiers de longueur binaire $O(|X| +
        \log|V| + 1/\epsilon)$, telle qu' un ensemble de colonnes de $M$ est indépendant avec une
        probabilité $\epsilon$ si $T$ est lié à $S$ dans $G - R$.
    \end{rmq}
\end{frame}
%}}}
\section{La kernalisation}
\subsection*{a}

%{{{ Kernalisation
\begin{frame}
    Soient une instance $(G, k)$ de \oct et un paramètre d'erreur $\epsilon$ :
    \vfill
    \begin{enumerate}
        \item Si $k < \log(n)$, exécuter l'algorithme de Reed, Smith et Vetta et retourner une OUI
            ou NON instance de taille constante
            \pause
        \item Sinon, solution approchée $X$ de ratio $O(\sqrt{\log n}) = O(k^{1/2})$
            \pause
        \item Si $|X| \leq k$ retourner une OUI instance
            \pause
        \item Sinon construire $G'$ à partir de $G$ et de $X$
            \pause
        \item Appliquer l'algorithme randomisé pour obtenir une matrice
            \pause
        \item $(A, k)$ est une compression polynomiale de $(G, k)$
    \end{enumerate}
\end{frame}
% }}}

\section{Conclusion}
\subsection*{a}

\begin{frame}
\end{frame}

\end{document}
